\section{Discussion}
\label{sec:discussion}

\paragraph{A protocol, not another motif count.}
The zero-common-relay counterexample isolates a distinction that motif
counts cannot express. Separate routes expose both sources in one
reconstruction, while the other loses a source at the readout. This
source-access difference is consistent with the archived high/chance
outcomes and complements boundary-aware controls~\cite{park2026}. A
compensated deletion must therefore be interpreted together with its
replacement routes.

\paragraph{What the criterion guarantees.}
Proposition~1 is a population statement for one noiseless subspace task.
It says whether the XOR label is decodable from transferred source states;
it does not promise finite-MLP accuracy or biological generality.

\paragraph{Finite trainability is separate.}
All nonzero $k=8$ doses preserve ideal label information, while finite
accuracy varies substantially. Those curves describe how the recorded
procedure handles different input mixtures within an information-preserving
regime. Rank-one transfer can retain the label, so full source recovery is a
stronger requirement than this classification task needs.

\paragraph{Scope.}
The archived task uses synthetic labels, two-step linear propagation, and a
finite MLP on subgraphs of one male connectome. The $k=8$ extractions share
195 neurons. The CIFAR-10 assay is a readout test on four mixed image codes,
not a recognition benchmark. The population result assumes noiseless
continuous features; nonlinear propagation requires input-dependent
analysis. The reverse-support CIFAR subset detects no connectome-specific
benefit. The two dense $k=8$ bases match frozen source hashes, but original
dose-instance hashes are unavailable. Reconstructed operators give a
consistency result, and the edit algorithms differ in order and partner
choices, so the historical gap cannot be attributed to one design rule.

\section{Conclusion}
\label{sec:conclusion}

For the fixed connectome propagator studied here, source visibility reports
the task information that survives an intervention. On the noiseless subspace
task, it yields an exact Bayes criterion. On matched narrow graphs, the
criterion separates zero-common-relay edits with different source access and
distinguishes ideal information from finite-readout dose curves. Executed
visual readouts further locate what the audit does and does not cover. The
broader reporting lesson is direct: a
preserved structural statistic certifies only the constraint it enforces.
Before interpreting an intervention, state the executed operator and test
what information reaches its readout.
